\section{Comercio internacional, aranceles, monopolio y fallas de mercado}
\label{sec:comercio-monopolio-fallas}

En esta parte se conectan varios temas fundamentales del curso: la lógica del comercio internacional, la intervención estatal a través de aranceles, el poder de mercado del monopolio y la existencia de fallas del mercado como externalidades y bienes públicos.

\subsection{Comercio internacional}

El comercio internacional ocurre porque los países tienen diferentes costos de oportunidad, recursos, tecnología o dotación de factores. Esto genera especialización y permite consumir más y mejor de lo que sería posible con autarquía.

\begin{concepto}{Autarquía}
La autarquía es una economía cerrada, que no comercia con el exterior.
\end{concepto}

\begin{itemize}
    \item \textbf{Exportaciones:} venta de bienes al exterior.
    \item \textbf{Importaciones:} compra de bienes desde el exterior.
\end{itemize}

\begin{ejemplo}
Si Chile compara muy bien en vino y Argentina en trigo, ambos ganan si Chile exporta vino e importa trigo, mientras Argentina exporta trigo e importa vino.
\end{ejemplo}

\subsection{Ganancias del comercio}

\begin{concepto}{Ganancias del comercio}
El comercio aumenta el bienestar porque permite especialización, menores costos de producción y mayor variedad de bienes.
\end{concepto}

Cuando un país se abre al comercio, puede consumir afuera de su frontera de posibilidades de producción original.

\begin{ejemplo}
Un país sin comercio puede consumir solo lo que produce internamente. Con comercio puede consumir una combinación que antes era imposible, porque puede intercambiar producción especializada por otros bienes.
\end{ejemplo}

\subsection{Importaciones y exportaciones}

\begin{itemize}
    \item Si el precio internacional es menor que el precio interno, el país importará.
    \item Si el precio internacional es mayor que el precio interno, el país exportará.
\end{itemize}

\begin{observacion}{Importante}
Importar puede beneficiar a los consumidores, pero perjudica a los productores locales que compiten con bienes importados.
\end{observacion}

\subsection{Aranceles}

Un arancel es un impuesto que grava la importación de bienes. Su propósito puede ser proteger a la industria local, pero implica costos para la sociedad.

\begin{concepto}{Arancel}
Impuesto que se aplica a un bien importado.
\end{concepto}

\begin{ejemplo}
Si el precio mundial de un bien es menor que el precio interno, un arancel hace subir el precio interno, reduciendo las importaciones y aumentando la producción local.
\end{ejemplo}

\subsection{Efectos de un arancel}

\begin{itemize}
    \item aumenta el precio interno,
    \item disminuye la cantidad demandada,
    \item aumenta la producción local,
    \item reduce importaciones,
    \item genera recaudación fiscal para el Estado,
    \item pero crea pérdida social.
\end{itemize}

\begin{advertencia}{Trampa de prueba}
No basta con decir que un arancel “protege la industria local”. También debe analizarse el costo para los consumidores y la pérdida de bienestar total.
\end{advertencia}

\subsection{Pérdida social del arancel}

El arancel genera dos tipos de ineficiencia:
\begin{enumerate}
    \item 	extbf{ineficiencia productiva:} se produce localmente algo que se podría importar a menor costo,
    \item 	extbf{subconsumo:} consumidores que valoran el bien por sobre el precio mundial dejan de comprarlo.
\end{enumerate}

\begin{ejemplo}
Si un país importa acero barato del exterior, un arancel puede provocar que se produzca acero más caro localmente, aunque la sociedad tendría menos costo si importara ese acero.
\end{ejemplo}

\subsection{Monopolio}

Un monopolio es una estructura de mercado en la que hay un solo vendedor de un bien o servicio.

\begin{concepto}{Monopolio}
Es una situación donde una empresa tiene poder de mercado suficiente para fijar el precio y la cantidad.
\end{concepto}

\subsection{Causas del monopolio}

Los monopolios surgen por diferentes razones:
\begin{itemize}
    	\item patentes,
    \item exclusividad de un recurso,
    \item barreras legales o regulatorias,
    \item economías de escala,
    \item control de infraestructura esencial.
\end{itemize}

\begin{ejemplo}
Una empresa que posee una patente para un medicamento puede convertirse en monopolio temporal, porque nadie más puede producir ese medicamento sin permiso.
\end{ejemplo}

\subsection{Monopolio y pérdida social}

El monopolio produce una cantidad menor y un precio mayor que la competencia perfecta. Esto genera pérdida social porque se dejan de realizar intercambios mutuamente beneficiosos.

\begin{observacion}{Importante}
Un monopolio no es necesariamente “malo” por definición, pero sí suele generar menor eficiencia que mercados competitivos.
\end{observacion}

\subsection{Externalidades}

Una externalidad es un efecto que una actividad tiene sobre terceros y que no está incorporado en el precio de mercado.

\begin{concepto}{Externalidad}
Es un costo o beneficio no internalizado por el agente que tomó la decisión.
\end{concepto}

\subsubsection{Externalidades negativas}

Son costos que recaen sobre terceros sin compensación.

\begin{ejemplo}
La contaminación de una fábrica afecta a los vecinos y a la salud pública, pero no siempre el fabricante paga ese costo completo.
\end{ejemplo}

\subsubsection{Externalidades positivas}

Son beneficios que reciben terceros sin haber pagado por ellos.

\begin{ejemplo}
La vacunación reduce la propagación de enfermedades para toda la comunidad, no solo para quien se vacuna.
\end{ejemplo}

\subsection{Fallas de mercado y soluciones}

Cuando hay externalidades, el mercado puede producir demasiado o demasiado poco de ciertos bienes.

\begin{itemize}
    \item Externalidad negativa: suele producirse demasiado.
    \item Externalidad positiva: suele producirse demasiado poco.
\end{itemize}

Las soluciones pueden ser:
\begin{itemize}
    \item impuestos para bienes con externalidades negativas,
    \item subsidios para bienes con externalidades positivas,
    \item regulación,
    \item derechos de propiedad,
    \item acuerdos privados (teorema de Coase).
\end{itemize}

\subsection{Bienes públicos y recursos comunes}

\begin{itemize}
    \item \textbf{Bienes públicos:} no rivales y no excluibles.
    \item \textbf{Recursos comunes:} rivales pero no excluibles.
    \item \textbf{Bienes privados:} rivales y excluibles.
    \item \textbf{Bienes club:} no rivales y excluibles.
\end{itemize}

\begin{ejemplo}
- Bien público: faro, defensa nacional.\
- Recurso común: peces del mar, bosques en riesgo.\
- Bien privado: una manzana.\
- Bien club: televisión por suscripción.
\end{ejemplo}

\subsection{Importancia de esto en la economía}

Las fallas de mercado muestran que el mercado no siempre resuelve eficientemente todos los problemas sociales. Por eso, la intervención pública puede ser necesaria para corregir distorsiones, aunque no siempre es la solución más eficiente si no se diseñan bien las políticas.

\begin{resumen}{Puntos clave}
- El comercio genera ganancias por especialización.
- Los aranceles protegen sectores, pero reducen bienestar agregado.
- El monopolio genera precio alto y cantidad baja.
- Las externalidades justifican intervención pública.
- Los bienes públicos y recursos comunes requieren regulación especial.
\end{resumen}
